\subsection{扭曲线丛与 Poincar\'e 丛}\label{src:9.4.4-subsection}
称 $\cBun_{r,d}$ 上线丛 $L$ 为 $w$-扭曲的, 若对每个向量丛 $F$, 数乘自同构 $\mathbb G_m\subseteq\underline{\operatorname{Aut}}(F)$ 在 $L\otimes\kappa([F])$ 上的作用权重为 $w$. 因为 $\cBun_{r,d}$ 连通, 权重 $w$ 与 $F$ 的选取无关. 例如, 行列式线丛 $\mathcal L_V$ 是 $\chi(C,F\otimes V)$-扭曲的, 因为 $\mathcal L_V\otimes\kappa([F])\cong\det\mathrm R\Gamma(C,F\otimes V)$.
\begin{proposition}\label{src:9.4.30}
Picard 叠 $\underline{\mathcal P\!ic}^d$ 有 $1$-扭曲线丛, 且带 $\mathbb G_m$-带的 gerbe $\underline{\mathcal P\!ic}^d\to\mathrm{Pic}^d$ 平凡.
\end{proposition}
\begin{proof}
若 $V$ 的秩为 $1$, 次数为 $g-d$, 则对 $L\in\mathrm{Pic}^d(C)$, 有 $\chi(C,L\otimes V)=\deg(L\otimes V)+(1-g)=1$. 所以 $\mathcal L_V$ 是 $1$-扭曲的. 第二个断言由习题~\bookref{7.4.49} 得到: $\mathbb G_m$-gerbe 平凡, 当且仅当有 $1$-扭曲线丛.
\end{proof}
更一般地, 可刻画 $\cBun_{r,d}^{\mathrm{ss}}$ 上扭曲线丛的可能权重.
\begin{proposition}\label{src:9.4.31}
$\cBun_{r,d}^{\mathrm{ss}}$ 上存在 $w$-扭曲线丛, 当且仅当 $w$ 为 $\gcd(r,d)$ 的倍数.
\end{proposition}
\begin{proof}
写为 $k:=\gcd(r,d)=ar+bd$, 其中 $a,b$ 为整数. 若 $V$ 是秩 $b$, 次数 $a-b(1-g)$ 的向量丛, 则断言 $\mathcal L_V$ 为 $k$-扭曲的, 因而对任何整数 $m$, $\mathcal L_V^{\otimes m}$ 为 $mk$-扭曲的. 为此, 对秩 $r$, 次数 $d$ 的向量丛 $F$ 直接计算
\[
\chi(C,F\otimes V)=db+r(a-b(1-g))+rb(1-g)=k.
\]
反之, 若 $L$ 是 $\cBun_{r,d}^{\mathrm{ss}}$ 上 $w$-扭曲线丛, 则权重 $w$ 可计算为在 $G^{\oplus k}$ 处的特征标, 其中 $G$ 是秩 $r':=r/k$, 次数 $d':=d/k$ 的稳定向量丛. 在由 $E\mapsto E^{\oplus k}$ 定义的映射 $f:\cBun_{r',d'}\to\cBun_{r,d}$ 下, $L$ 在 $G^{\oplus k}$ 处的特征标是 $f^*L$ 在 $G$ 处的特征标的 $k$ 倍; 后者必为整数. 所以 $k$ 整除 $w$.
\end{proof}
\begin{definition}\label{src:9.4.32}
$C\times M_{r,d}^{\mathrm{ps}}$ 上的\newterm{Poincar\'e 丛}{Poincar\'e bundle} (或万有向量丛\footnote{我们更倾向于使用``Poincar\'e 丛'', 以免与模叠 $C\times\cBun_{r,d}^{\mathrm{ss}}$ 上的万有向量丛 $\mathcal F_{\mathrm{univ}}$ 混淆.}) 是秩 $r$, 次数 $d$ 的半稳定向量丛 $E_{\mathrm{Poin}}$, 使得对每个对应于 $C\times T$ 上半稳定层 $F$ 的映射 $f:T\to\cBun_{r,d}^{\mathrm{ss}}$, 存在线丛 $L\in\mathrm{Pic}(T)$ 和 $C\times T$ 上向量丛的同构
\[
F\cong(\mathrm{id}\times(\pi\circ f))^*E_{\mathrm{Poin}}\otimes p_2^*L,
\]
其中 $\pi$ 如常表示 $\cBun_{r,d}^{\mathrm{ss}}\to M_{r,d}^{\mathrm{ps}}$.
\end{definition}
\begin{remark}\label{src:9.4.33}
换言之, $E_{\mathrm{Poin}}$ 是万有丛, 若存在 $\cBun_{r,d}^{\mathrm{ss}}$ 上的线丛 $L$, 使得 $\mathcal F_{\mathrm{univ}}\cong(\mathrm{id}\times\pi)^*E_{\mathrm{Poin}}\otimes p_2^*L$ 为 $C\times\cBun_{r,d}^{\mathrm{ss}}$ 上向量丛的同构. 等价地, 万有丛 (若存在) 是函子 $T\mapsto\cBun_{r,d}^{\mathrm{ss}}(T)/p_2^*\mathrm{Pic}(T)$ 的万有族.
\end{remark}
\begin{theorem}\label{src:9.4.34}
限制到稳定轨迹 $\cBun_{r,d}^{\mathrm s}\to M_{r,d}^{\mathrm s}$ 后, 下列条件等价:
\begin{enumerate}[(1)]
\item $C\times M_{r,d}^{\mathrm s}$ 上存在 Poincar\'e 丛 $E_{\mathrm{Poin}}$;
\item $\mathbb G_m$-gerbe $\cBun_{r,d}^{\mathrm s}\to M_{r,d}^{\mathrm s}$ 平凡;
\item $\cBun_{r,d}^{\mathrm s}$ 上存在 $1$-扭曲线丛.
\end{enumerate}
\end{theorem}
\begin{proof}
为说明 (1) $\Leftrightarrow$ (2), 注意 $C\times M_{r,d}^{\mathrm{ps}}$ 上任何半稳定向量丛 $E$ 都诱导映射 $\sigma:\cBun_{r,d}^{\mathrm s}\to M_{r,d}^{\mathrm s}$, 可以验证复合 $\pi\circ\sigma$ 为恒等映射 (即 $\sigma$ 为截面), 当且仅当 $E$ 为 Poincar\'e 丛. 等价性 (2) $\Leftrightarrow$ (3) 来自习题~\bookref{7.4.49}.
\end{proof}
\begin{corollary}\label{src:9.4.35}
若 $\gcd(r,d)=1$, 则 $\cBun_{r,d}^{\mathrm{ss}}\to M_{r,d}^{\mathrm{ps}}$ 是平凡 $\mathbb G_m$-gerbe, 并且 $C\times M_{r,d}^{\mathrm{ps}}$ 上存在 Poincar\'e 丛.
\end{corollary}
\begin{proof}
由 $1$-扭曲线丛的存在性 (命题~\bookref{9.4.31}) 以及定理~\bookref{9.4.34} 得到. 另见 \cite[定理 8.4.2]{LP97}.
\end{proof}
